\chapter{Marco teórico}
\label{chap:marco}

Este capítulo resume los conceptos necesarios para seguir el resto de la memoria: cómo se acciona
el eje, cómo se mide la fuerza, con qué microcontrolador y por qué buses viaja la información.

\section{Transmisión por cable y radio variable}

En una transmisión por cable, el par que hay que aplicar al tambor es el producto de la tensión del
cable por el radio con el que ese cable sale del tambor:
\begin{equation}
  \tau = T\,r .
  \label{eq:par-basico}
\end{equation}
Si el radio es constante, el par reproduce exactamente la forma de la tensión. Si el radio cambia
con el giro, $r(\varphi)$, se dispone de un grado de libertad para moldear el par: basta con que el
radio sea grande donde la tensión es pequeña, y pequeño donde es grande, para que el producto se
mantenga aproximadamente constante. Ésa es la idea del \emph{fusee} y de los tambores de radio
variable, y el capítulo~\ref{chap:modelo} la desarrolla para la geometría concreta del prototipo.

Conviene separar dos cosas que se confunden con facilidad. Un mecanismo de compensación
\emph{pasiva} (un muelle con una polea no circular) reduce la energía que el actuador tiene que
aportar. Una caracola accionada por un motor no reduce esa energía: el trabajo de subir la masa es
el mismo. Lo que hace es repartirlo de manera uniforme a lo largo del giro, lo que reduce el par de
pico y, con él, el tamaño del motor necesario.

\section{Accionamiento: servomotores paso a paso de lazo cerrado}

Un motor paso a paso convencional funciona en lazo abierto: el controlador impone la secuencia de
fases y se da por hecho que el rotor la sigue. Es sencillo y tiene buen par a baja velocidad, pero
si la carga supera el par disponible se pierden pasos sin que el sistema lo detecte
\cite{kenjo1994}. Los servomotores paso a paso de lazo cerrado añaden un codificador y un
accionamiento que regula la corriente en función del error de posición: conservan el par a baja
velocidad, no pierden pasos y ofrecen modos de control en posición, velocidad o par, además de
magnitudes de supervisión como la velocidad real o el par entregado.

El accionamiento de este trabajo pertenece a esa familia \completar{marca y modelo del motor y del
accionamiento, par nominal, tensión de alimentación y resolución del codificador}. Integra una
interfaz RS-485 con protocolo Modbus RTU, por la que se configura (modo de control, rigidez, rampas
y resistencia de frenado), se habilita, se envía la consigna de velocidad y se leen la velocidad y
el par. Tenerlo todo en un bus de dos hilos simplifica el cableado frente a la alternativa de
señales de paso y dirección más entradas analógicas.

\section{Medida de fuerza con células de carga}

\subsection{Galgas extensométricas y puente de Wheatstone}

Una célula de carga convierte la fuerza en una deformación elástica, que se mide con galgas
extensométricas dispuestas en un puente de Wheatstone \cite{hoffmann2012}. La salida del puente es
una tensión diferencial de unos pocos milivoltios por voltio de excitación a plena escala, que hay
que amplificar con alta ganancia y bajo ruido antes de digitalizarla; además hay que compensar el
desequilibrio de cero del puente y su deriva con la temperatura \cite{pallas2001}. De ahí que la
cadena de medida necesite, ante todo, una calibración con masas patrón: sin ella se conoce la forma
de la señal, pero no su escala ni su cero.

\subsection{Acondicionadores integrados}

Existen circuitos integrados que agrupan excitación, amplificación, conversión A/D y corrección
digital. En este trabajo se han evaluado dos:

\begin{itemize}
  \item \textbf{ZSC31014} de Renesas \cite{zsc31014}, con salida digital I\textsuperscript{2}C,
        preamplificador de ganancia programable, conversor de 14 bits y corrección almacenada en
        EEPROM. Cada lectura incluye dos bits de estado que indican si el dato es nuevo. Su
        configuración sólo puede modificarse en un \emph{modo comando} al que se accede durante una
        ventana corta tras el encendido, lo que ha condicionado el diseño
        (sección~\ref{sec:prob-zsc}). Se usa en una placa \emph{Load Cell 2 Click}
        \cite{mikroelc2}.
  \item \textbf{NAU7802} de Nuvoton \cite{nau7802}, conversor sigma-delta de 24 bits con
        amplificador programable y regulador interno. Se probó como segunda célula y se retiró del
        firmware final.
\end{itemize}

\section{Microcontrolador}

El proyecto empezó con una placa basada en ESP32-S3 (Seeed XIAO) y se migró a una Raspberry Pi
Pico~2. La tabla~\ref{tab:mcu} resume las características del RP2350 \cite{rp2350}. El motivo del
cambio y sus consecuencias se explican en la sección~\ref{sec:prob-uart}: los pines que impone el
trazado de la placa no admiten una UART hardware, pero el RP2350 dispone de periféricos PIO capaces
de implementar una UART en cualquier pin.

\begin{table}[H]
\centering
\caption{Características del microcontrolador RP2350 en la Raspberry Pi Pico 2 \cite{rp2350}.}
\label{tab:mcu}
\begin{tabular}{@{}ll@{}}
\toprule
Característica & Raspberry Pi Pico 2 (RP2350) \\
\midrule
Núcleos & 2 × Arm Cortex-M33 o 2 × RISC-V Hazard3, a \SI{150}{\mega\hertz} \\
Memoria & \SI{520}{\kilo\byte} SRAM, \SI{4}{\mega\byte} flash QSPI \\
Periféricos serie & 2 × UART, 2 × SPI, 2 × I\textsuperscript{2}C \\
E/S programable & 12 máquinas de estado PIO \\
Entradas analógicas & 3 canales ADC \\
Entorno usado & Arduino-Pico \cite{arduinopico} \\
\bottomrule
\end{tabular}
\end{table}

\section{Buses de comunicación}

\subsection{RS-485 y Modbus RTU}

RS-485 \cite{tia485} es un estándar de capa física diferencial y multipunto, extendido en la
industria por su inmunidad al ruido, su alcance y la posibilidad de conectar varios equipos al mismo
par trenzado. En modo semidúplex cada nodo activa su emisor sólo mientras transmite, y ese control
de dirección (líneas DE y $\overline{\text{RE}}$ del transceptor) corresponde al maestro.

Modbus \cite{modbusapp} es un protocolo maestro-esclavo de petición y respuesta que modela los datos
del esclavo como registros de 16 bits. En su variante RTU sobre línea serie \cite{modbusserial},
cada trama lleva la dirección del esclavo, el código de función, los datos y un CRC-16, y las tramas
se delimitan por silencios de al menos 3,5 caracteres. Aquí se usan la función 0x03 (lectura de
registros) y la 0x06 (escritura de un registro), a través de la biblioteca ModbusMaster
\cite{modbusmaster}.

\subsection{I\textsuperscript{2}C y mikroBUS}

I\textsuperscript{2}C \cite{i2cspec} es un bus síncrono de dos hilos con resistencias de
\emph{pull-up}, pensado para comunicar circuitos dentro de una misma placa; es el que usan los
acondicionadores de célula de carga. Las placas Click siguen el estándar mikroBUS \cite{mikrobus},
que define un zócalo de 16 pines con I\textsuperscript{2}C, SPI, UART, alimentación y señales
auxiliares.

\section{Adquisición y supervisión en el ordenador}

La supervisión de máquinas se ha hecho tradicionalmente con paneles HMI o paquetes SCADA
comerciales. Para un banco de laboratorio, una aplicación web local es una alternativa ligera: no
requiere licencias, funciona en cualquier navegador y admite varios clientes a la vez. En este
trabajo se emplea FastAPI \cite{fastapi} con pySerial \cite{pyserial} para el puerto USB, WebSocket
\cite{rfc6455} para el envío continuo de datos al navegador y Chart.js \cite{chartjs} para las
gráficas.

Dos conceptos de adquisición son importantes para interpretar los resultados. El primero es la
\emph{base de tiempos}: si las muestras se marcan con el reloj del ordenador, la latencia del enlace
se confunde con el fenómeno medido; por eso aquí cada muestra lleva la marca de tiempo del
microcontrolador. El segundo es la distinción entre \emph{dato bruto} y \emph{dato calibrado}:
guardar las cuentas del conversor, y no los newtons, permite recalibrar un ensayo sin repetirlo,
que es justo lo que ha hecho falta en este trabajo.
